\section{Kritéria výběru}
\label{sec:selection-criteria}

Aby bylo možné objektivně porovnat dostupné modely a způsoby nasazení, je třeba nejprve definovat kritéria, podle nichž bude hodnocení probíhat.
Kritéria vycházejí z požadavků zadání diplomové práce, z technických omezení systému Kelvin a z praktických podmínek školního prostředí.

\subsection{Kvalita výstupů}
\label{subsec:criteria-quality}

Primárním účelem nasazeného modelu je generovat komentáře ke studentskému zdrojovému kódu, které jsou věcně správné, relevantní ke konkrétnímu řešení a formulované tak, aby je vyučující mohl přímo využít jako základ zpětné vazby.
Kvalita a relevance výstupu je proto nejdůležitějším kritériem.

Z hlediska správnosti je klíčové, aby model identifikoval skutečné problémy v kódu a nenahlašoval falešné poplachy.
Falešně pozitivní komentáře vyučujícího zdržují a snižují důvěru v celý systém.
Stejně tak jsou problematické příliš obecné připomínky, které nenesou konkrétní informaci o tom, co a kde v kódu je problematické.

Relevance výstupů úzce závisí na schopnosti modelu pracovat s kontextem zadání.
Model, který vidí pouze zdrojový kód bez popisu úlohy, může identifikovat obecné stylové nedostatky, ale nedokáže posoudit, zda zvolená datová struktura odpovídá požadavkům zadání, nebo zda výstup programu splňuje specifikaci.

\subsection{Schopnost pojmout celé studentské řešení}
\label{subsec:criteria-context}

Jak bylo vysvětleno v sekci~\ref{subsec:llm-tokens-parameters-context}, každý model má omezené kontextové okno dané maximálním počtem tokenů, které zpracuje v jednom průchodu.
Pro analýzu studentských řešení je nutné, aby model v tomto limitu zvládl:

\begin{itemize}
    \item systémový prompt specifikující chování modelu a výstupní formát,
    \item případné text zadání úlohy,
    \item kompletní zdrojový kód odevzdaného řešení,
    \item dostatečnou rezervu pro vygenerování odpovědi.
\end{itemize}

Typické studentské řešení domácí úlohy v jazyce C++ se pohybuje v rozsahu 100 až 500 řádků kódu.
Po tokenizaci a s přidáním zadání a systémového promptu se celkový vstup obvykle vejde do 3\,000 až 5\,000 tokenů.
Modely s kontextovým oknem alespoň 8\,000 tokenů jsou proto pro tento účel dostačující.
Modely s výrazně menším oknem by vyžadovaly dělení vstupu na části, což komplikuje implementaci a snižuje kvalitu analýzy, protože model nemá přístup k celku řešení najednou.

V případě odevzdání rozsáhlejšího projektu nebo sady souborů je nicméně vhodné mít k dispozici model s kontextem v řádu desítek tisíc tokenů, aby bylo možné tyto situace zvládnout bez nutnosti chunking strategie.

\subsection{Dostupnost a udržovatelnost}
\label{subsec:criteria-availability}

Školní informační systém musí fungovat spolehlivě po dobu celého semestru, přičemž nejvyšší zátěž nastává před deadlinem odevzdání úloh.
Zvolený model nebo způsob přístupu k němu musí zaručit dostatečnou dostupnost a stabilitu.

U komerčních API je dostupnost garantována licencí a poskytovatelé mají typicky vysokou míru dostupnosti.
Klíčové jsou však i podmínky, zejména to, zda mohou být výstupy modelu použity v prostředí školního hodnocení, a zda se podmínky API v čase nemění způsobem, který by narušil provoz systému.

U open-source modelů nasazených lokálně závisí dostupnost na spolehlivosti vlastní infrastruktury.
Výhodou je plná kontrola nad verzí modelu, model lze fixovat na konkrétní verzi a nevzniká riziko, že by poskytovatel model stáhl nebo výrazně změnil.

Z pohledu udržovatelnosti je důležité i to, jak snadné je model aktualizovat nebo vyměnit za novější verzi.
Integrace by měla být navržena tak, aby výměna modelu nevyžadovala žádné změny v kódu systému Kelvin.

\subsection{Náročnost provozu}
\label{subsec:criteria-operational}

Provozní náklady a náročnost se liší podle zvoleného způsobu nasazení.
Při cloudovém API jsou náklady přímou funkcí objemu zpracovaných tokenů -- za každý dotaz je účtován poplatek.
Pro odhad nákladů je proto nutné znát přibližný počet odevzdání za semestr a průměrnou délku zpracovávaného vstupu.


\endinput